\documentclass[11pt,a4paper]{article}

% Pakete für deutsche Sprache
\usepackage[ngerman]{babel}
\usepackage[utf8]{inputenc}
\usepackage[T1]{fontenc}

% Weitere wichtige Pakete
\usepackage{graphicx}
\usepackage{listings}
\usepackage{color}
\usepackage{hyperref}
\usepackage{geometry}
\usepackage{fancyhdr}
\usepackage{float}
\usepackage{caption}
\usepackage{subcaption}
\usepackage{setspace}

% Schriftart auf Arial (Helvetica) setzen
\usepackage[scaled]{helvet}
\renewcommand{\familydefault}{\sfdefault}

% Zeilenabstand 1,5-fach
\onehalfspacing

% Blocksatz mit Silbentrennung
\usepackage{microtype}
\sloppy

% Seitenränder entsprechend IHK-Vorgaben
\geometry{
    left=2.5cm,
    right=2.5cm,
    top=4cm,
    bottom=3cm,
    headheight=3cm,
    headsep=0.5cm,
    footskip=1cm
}

% Code-Highlighting Setup
\definecolor{codegreen}{rgb}{0,0.6,0}
\definecolor{codegray}{rgb}{0.5,0.5,0.5}
\definecolor{codepurple}{rgb}{0.58,0,0.82}
\definecolor{backcolour}{rgb}{0.95,0.95,0.92}

\lstdefinestyle{mystyle}{
    backgroundcolor=\color{backcolour},   
    commentstyle=\color{codegreen},
    keywordstyle=\color{magenta},
    numberstyle=\tiny\color{codegray},
    stringstyle=\color{codepurple},
    basicstyle=\ttfamily\footnotesize,
    breakatwhitespace=false,         
    breaklines=true,                 
    captionpos=b,                    
    keepspaces=true,                 
    numbers=left,                    
    numbersep=5pt,                  
    showspaces=false,                
    showstringspaces=false,
    showtabs=false,                  
    tabsize=2
}

\lstset{style=mystyle}

% Kopf- und Fußzeile
\pagestyle{fancy}
\fancyhf{}
\renewcommand{\headrulewidth}{0pt}
\renewcommand{\footrulewidth}{0pt}
\setlength{\headheight}{3cm}
\setlength{\headsep}{0.5cm}
\setlength{\footskip}{1cm}
\lhead{\hspace*{-2.0cm}\vspace*{2cm}\includegraphics[height=2.2cm]{image1.png}}
\cfoot{\thepage}

\begin{document}

% Titelseite
\thispagestyle{empty}

% Logo-Bereich oben
\begin{figure}[H]
    \hspace*{-2.0cm}\vspace*{2cm}\includegraphics[height=2.2cm]{image1.png}
\end{figure}

\vspace{3cm}

% Haupttitel - zentriert
\begin{center}
    {\Huge \textbf{COSPAR Communication System}}
\end{center}

\vfill

% Informationsblock unten
\begin{center}
    \begin{tabular}{ll}
        \textbf{Prüfling:} & Abdullah Diab \\[0.5em]
        \textbf{Geburtsdatum:} & 01.01.2006 \\[0.5em]
        \textbf{Ausbildungsberuf:} & Fachinformatiker für Anwendungsentwicklung \\[0.5em]
        \textbf{Ausbildungsbetrieb:} & AFZ - Aus- und Fortbildungszentrum Bremen \\[0.5em]
        \textbf{Ausbildungsstelle:} & ZARM - Zentrum für angewandte Raumfahrttechnologie und Mikrogravitation \\[0.5em]
        \textbf{Prüflings-Nr.:} & 40016 \\
    \end{tabular}
\end{center}

\vspace{2cm}

\newpage

% Eidesstattliche Erklärung
\section*{Eidesstattliche Erklärung}
\thispagestyle{empty}

Hiermit versichere ich, dass das Projekt mit dem Titel \glqq COSPAR Communication System\grqq{} und die dazugehörige Dokumentation in Einzelarbeit von meiner Person angefertigt wurde. Diese Dokumentation wurde in der Vergangenheit nicht bei Prüfungen zur Bewertung/Begutachtung vorgelegt.

\vspace{3cm}

\noindent
Bremen, den \underline{\hspace{3cm}}

\vspace{2cm}

\noindent
\underline{\hspace{8cm}}\\
Ort, Datum und Unterschrift des Prüfungsteilnehmers

\newpage

% Inhaltsverzeichnis
\pagenumbering{roman}
\setcounter{page}{1}
\tableofcontents
\newpage

% Abbildungsverzeichnis
\thispagestyle{empty}
\listoffigures

% Ab hier beginnt die eigentliche Dokumentation mit arabischen Seitenzahlen
\clearpage
\pagenumbering{arabic}
\setcounter{page}{1}
\section{Einleitung}

\subsection{Der Ausbildungsbetrieb}
% Beschreibung des ZARM als Ausbildungsbetrieb
Das Zentrum für angewandte Raumfahrttechnologie und Mikrogravitation (ZARM) der Universität Bremen ist eine international anerkannte Forschungseinrichtung, die sich auf Grundlagenforschung und angewandte Forschung in der Raumfahrttechnologie und Mikrogravitationsforschung spezialisiert hat. Als interdisziplinäres Zentrum verbindet das ZARM Expertise aus den Bereichen Physik, Ingenieurswissenschaften und Materialwissenschaften.

% Rolle im internationalen Wissenschaftsbetrieb
Das ZARM spielt eine bedeutende Rolle in der internationalen Raumfahrtforschung und unterhält enge Kooperationen mit führenden Raumfahrtagenturen wie der ESA, NASA und DLR. Durch diese Vernetzung ist das ZARM auch aktiv in wissenschaftlichen Organisationen wie COSPAR (Committee on Space Research) vertreten, wo Mitarbeiter wichtige Funktionen als Main Scientific Organizers (MSOs) und Deputy Organizers (DOs) übernehmen.

% Ausbildungskontext  
Im Rahmen der Ausbildung werden praktische IT-Projekte durchgeführt, die reale Problemstellungen aus dem wissenschaftlichen Betrieb aufgreifen und digitale Lösungen entwickeln. Das COSPAR Communication System Projekt ist ein solches Ausbildungsprojekt, das eine konkrete Herausforderung im COSPAR-Organisationsprozess adressiert.

\subsection{Zielsetzung des Projekts}
Das COSPAR Communication System Projekt verfolgt mehrere zentrale Ziele zur Verbesserung der Kommunikationsprozesse zwischen MSOs/DOs und ihren Session-Autoren. Ein Hauptziel ist die deutliche Steigerung der Effizienz, sodass MSOs und DOs ihre Autoren wesentlich schneller kontaktieren können als bisher. Dies wird durch eine Bulk-E-Mail-Funktionalität erreicht, die es ermöglicht, mehrere Autoren gleichzeitig zu kontaktieren.

Ein weiterer wichtiger Aspekt ist die Entwicklung intelligenter Filteroptionen, die eine präzise Auswahl nach verschiedenen Kriterien ermöglichen. Diese umfassen die Filterung nach Sessions, Autorentypen wie Presenting Authors und Co-Authors sowie nach dem Präsentationsstatus. Dadurch können MSOs und DOs gezielt bestimmte Autorengruppen ansprechen.

Das Projekt eliminiert zeitaufwändige manuelle Prozesse durch automatisierte Systeme, die direkt auf die COSPAR-Datenbank zugreifen. Die nahtlose Integration in die bestehende COSPAR-Webinfrastruktur gewährleistet, dass das neue System harmonisch mit den vorhandenen Arbeitsabläufen funktioniert, ohne zusätzliche Komplexität für die Benutzer zu verursachen.

\subsection{Was ist COSPAR?}
COSPAR (Committee on Space Research) ist eine internationale wissenschaftliche Organisation, die 1958 vom International Council for Science (ICSU) gegründet wurde und als weltweites Forum für den wissenschaftlichen Austausch in der Weltraumforschung dient. Die Organisation koordiniert internationale Bemühungen in allen Bereichen der Raumfahrtwissenschaften und fördert die Zusammenarbeit zwischen verschiedenen Raumfahrtagenturen und Forschungseinrichtungen weltweit. Das wichtigste Ereignis von COSPAR sind die alle zwei Jahre stattfindenden Scientific Assemblies (Wissenschaftliche Versammlungen), bei denen Tausende von Wissenschaftlern aus aller Welt zusammenkommen, um aktuelle Forschungsergebnisse zu präsentieren und zu diskutieren. Diese Konferenzen sind nach wissenschaftlichen Kommissionen und Panels strukturiert, die verschiedene Bereiche der Weltraumforschung abdecken - von Planetenforschung über Astrophysik bis hin zu Weltraumtechnologie. COSPAR entwickelt wichtige Standards für Planetenschutz, koordiniert Datenformat-Richtlinien und organisiert den strukturierten Prozess der wissenschaftlichen Beitragsverwaltung über webbasierte Systeme. Die komplexe Organisation dieser internationalen Veranstaltungen erfordert ein effizientes Management-System, in dem die Main Scientific Organizers (MSOs) und Deputy Organizers (DOs) eine zentrale Rolle spielen, da sie die Schnittstelle zwischen der COSPAR-Führung und der wissenschaftlichen Gemeinschaft bilden.

\subsection{Die Rolle der MSOs, DOs und Autoren}
\textbf{Main Scientific Organizers (MSOs)}: Hauptverantwortliche für wissenschaftliche Sessions, die den gesamten Ablauf von der Abstract-Bewertung bis zur Session-Durchführung koordinieren. Sie müssen regelmäßig mit ihren Autoren kommunizieren, um Updates zu Präsentationen, Zeitplänen und technischen Anforderungen zu übermitteln.

\textbf{Deputy Organizers (DOs)}: Unterstützen MSOs bei der Session-Organisation und fungieren als zusätzliche Ansprechpartner. Sie benötigen ebenfalls direkten Zugang zu Autorenlisten für koordinierte Kommunikation.

\textbf{Autoren}: Wissenschaftler, die Abstracts einreichen und in verschiedenen Kategorien auftreten:
\begin{itemize}
    \item \textbf{Presenting Authors}: Hauptautoren, die ihre Forschung präsentieren
    \item \textbf{Co-Authors}: Mitautoren ohne Präsentationspflicht
\end{itemize}

\textbf{Kommunikationsbeziehungen}: MSOs/DOs müssen Autoren über Präsentationsdetails, Upload-Deadlines, technische Anforderungen und Zeitplanänderungen informieren. Bisher erfolgte dies manuell und zeitaufwändig.

\newpage

\section{Anforderungsanalyse}

\subsection{Ist-Analyse}

\subsubsection{Aktuelle Kommunikationsweise der MSOs/DOs}
Der bisherige Kommunikationsprozess erfolgt vollständig manuell und ist sehr zeitaufwändig. MSOs und DOs müssen zunächst ihr E-Mail-Programm wie Outlook oder Thunderbird öffnen und dann alle benötigten Autoren-E-Mail-Adressen einzeln aus der COSPAR-Datenbank kopieren. Für jede Kommunikation muss eine separate E-Mail erstellt werden, da kein systematisches Filterungssystem für verschiedene Autorengruppen existiert. Dieser Prozess wird besonders bei Sessions mit vielen Teilnehmern sehr zeitaufwändig und fehleranfällig.

\subsubsection{Probleme und Schwachstellen}
Die manuelle Vorgehensweise führt zu verschiedenen Problemen. Tippfehler bei der E-Mail-Adressen-Eingabe können dazu führen, dass wichtige Nachrichten nicht zugestellt werden. Bei umfangreichen Autorenlisten werden häufig einzelne Autoren vergessen, was zu unvollständiger Kommunikation führt. Die repetitiven Aufgaben bei jeder Kommunikationsrunde verschwenden wertvolle Arbeitszeit. Zusätzlich fehlt jede Art von Dokumentation oder Historienführung über gesendete E-Mails, und es gibt keine Möglichkeit zur gezielten Filterung nach Autorentypen oder Session-Eigenschaften.

\subsection{Soll-Analyse}

\subsubsection{Funktionale Anforderungen}
Das neue System soll eine sichere MSO/DO-Authentifizierung mit Berechtigungsprüfung für spezifische Sessions bieten. Kern des Systems sind session-basierte Autorenlisten, die automatisch alle Autoren einer bestimmten Session anzeigen. Erweiterte Filteroptionen ermöglichen die Auswahl nach Autorentyp (Presenting Authors oder Co-Authors), Präsentationstyp (Oral oder Poster Presentations) und Upload-Status der Präsentationen. Die Integration mit externen E-Mail-Clients erfolgt über automatische Mailto-Link-Generierung, während Such- und Filterfunktionen eine Echtzeit-Suche in Autorenlisten ermöglichen. Das responsive Design gewährleistet die Funktionalität auf verschiedenen Geräten.

\subsubsection{Nicht-funktionale Anforderungen}
Das System muss auch bei hunderten von Autoren pro Session eine schnelle Filterung gewährleisten und durch intuitive Bedienung maximal 3 Klicks zur E-Mail-Erstellung benötigen. Sichere Authentifizierung und datenschutzkonforme Handhabung von E-Mail-Adressen sind essentiell. Die browser-übergreifende Kompatibilität muss für Chrome, Firefox, Safari und Edge sichergestellt werden, während die nahtlose Integration in die bestehende COSPAR-Infrastruktur ohne Störung anderer Systeme erfolgen soll. Das System muss außerdem mit wachsenden Teilnehmerzahlen skalieren können.

\subsection{Wirtschaftlichkeitsbetrachtung}
Die Wirtschaftlichkeitsanalyse zeigt einen deutlichen Nutzen gegenüber dem Entwicklungsaufwand. Die größte Zeitersparnis liegt in der Reduktion der Kommunikationszeit von über 30 Minuten auf unter 2 Minuten pro Session. Gleichzeitig werden manuelle Eingabefehler bei E-Mail-Adressen eliminiert, was die Zuverlässigkeit der Kommunikation erheblich verbessert. Der Entwicklungsaufwand von etwa 35 Stunden führt zu langfristigen Effizienzgewinnen, wobei sich die Investition bereits nach wenigen COSPAR-Assemblies durch die eingesparte Arbeitszeit amortisiert.

\newpage
\section{Entwurfsphase}

\subsection{Systemarchitektur und Datenfluss}
Das COSPAR Communication System basiert auf einer modularen Architektur, die sich nahtlos in die bestehende COSPAR-Infrastruktur integriert:

\textbf{Datenfluss}:
\begin{enumerate}
    \item MSO/DO-Anmeldung über bestehendes COSPAR-Login-System
    \item Session-Identifikation und Berechtigungsprüfung
    \item Datenbankabfrage für Session-spezifische Autoren
    \item Anwendung von Filterkriterien (AJAX-basiert)
    \item Generierung von Mailto-Links für externe E-Mail-Clients
    \item Integration mit bestehenden E-Mail-Programmen der Benutzer
\end{enumerate}

\subsection{Technische Architektur}

\subsubsection{Analyse der bestehenden Codebasis}
Die vorhandene COSPAR-Webseite verwendet folgende Technologien:
\begin{itemize}
    \item \textbf{Backend}: PHP mit MySQL-Datenbank
    \item \textbf{Frontend}: HTML, CSS, JavaScript
    \item \textbf{Framework}: Eigenes MVC-ähnliches System
    \item \textbf{Authentifizierung}: Session-basierte Benutzeranmeldung
    \item \textbf{Datenbankzugriff}: mysqli-Verbindungen
\end{itemize}

\subsubsection{Wahl der Technologien}
Die Technologiewahl orientiert sich an der bestehenden COSPAR-Infrastruktur:
\begin{itemize}
    \item \textbf{Frontend}: HTML5, CSS3, Vanilla JavaScript für maximale Kompatibilität
    \item \textbf{Backend}: PHP 7.4+ für Konsistenz mit bestehenden Systemen
    \item \textbf{Datenbank}: MySQL mit bestehenden COSPAR-Tabellen
    \item \textbf{AJAX}: Für dynamische Filterung ohne Seitenneuladung
    \item \textbf{Responsive Design}: CSS Grid und Flexbox für Gerätekompatibilität
\end{itemize}

\subsubsection{Anbindung an die bestehende Datenbank}
Integration erfolgt über bestehende Tabellen:
\begin{itemize}
    \item \textbf{user}: Benutzerinformationen und Authentifizierung
    \item \textbf{sessions}: Session-Details und MSO/DO-Zuordnungen
    \item \textbf{abstracts}: Abstract-Informationen und Autorenverknüpfungen
    \item \textbf{authors}: Autoreninformationen und E-Mail-Adressen
    \item \textbf{presentations}: Upload-Status und Präsentationstypen
\end{itemize}

\subsection{Konzeption der Features}

\subsubsection{Autorenfilterung und -auswahl}
Konzept der intelligenten Autorenfilterung:
\begin{itemize}
    \item \textbf{Session-basierte Filterung}: Automatische Anzeige nur relevanter Autoren (A0.1, A0.2, etc.)
    \item \textbf{Autorentyp-Filter}: Unterscheidung zwischen Presenting Authors und Co-Authors
    \item \textbf{Kombinierte Filterkriterien}: Mehrere Filter gleichzeitig anwendbar
    \item \textbf{Ergebnistabelle}: Übersichtliche Darstellung mit Autoreninformationen
    \item \textbf{Auswahl-Checkboxen}: Individuelle Empfängersteuerung möglich
\end{itemize}

\subsubsection{Externe E-Mail-Client-Integration}
Konzept der nahtlosen Integration mit bestehenden E-Mail-Programmen:
\begin{itemize}
    \item \textbf{Mailto-Link-Generierung}: Automatische Erstellung von Links für externe E-Mail-Programme
    \item \textbf{Empfängerübernahme}: Automatische Übernahme gefilterter E-Mail-Adressen
    \item \textbf{Client-Kompatibilität}: Funktioniert mit Outlook, Thunderbird, Apple Mail, etc.
    \item \textbf{Bulk-E-Mail-Support}: Mehrere Empfänger in BCC für Datenschutz
    \item \textbf{Benutzerfreundlichkeit}: Ein-Klick-Lösung für E-Mail-Erstellung
\end{itemize}

\subsubsection{Benutzeroberflächen-Design}
Moderne und intuitive Benutzeroberfläche:
\begin{itemize}
    \item \textbf{Session-Header}: Farbkodierte Darstellung der Session-Informationen
    \item \textbf{Autorenkarten}: Card-basierte Anzeige mit wichtigen Informationen
    \item \textbf{Echtzeit-Suche}: Sofortige Filterung bei Texteingabe
    \item \textbf{Filter-Wizard}: Mehrstufiger Assistent für komplexe Filterkriterien
    \item \textbf{Responsive Design}: Optimale Darstellung auf allen Geräten
\end{itemize}

\subsubsection{Integration in bestehende COSPAR-Infrastruktur}
Nahtlose Einbindung in das bestehende System:
\begin{itemize}
    \item \textbf{Embedded-Modus}: Integration als Komponente in bestehende Seiten
    \item \textbf{Session-Management}: Nutzung der vorhandenen Authentifizierung
    \item \textbf{Datenkonsistenz}: Direkter Zugriff auf COSPAR-Datenbank
    \item \textbf{Design-Konsistenz}: Anpassung an bestehende Optik und Bedienung
\end{itemize}

\subsection{UI/UX Design}

\subsubsection{Wireframes und Mockups}
Das Design orientiert sich an modernen Web-Standards:
\begin{itemize}
    \item \textbf{Card-basierte Layouts}: Übersichtliche Darstellung von Autoreninformationen
    \item \textbf{Modal-Dialogs}: Filter-Wizard als überlagertes Fenster
    \item \textbf{Farbkodierung}: Session-spezifische Farben für bessere Orientierung
    \item \textbf{Responsive Grid}: Automatische Anpassung an verschiedene Bildschirmgrößen
\end{itemize}

\subsubsection{Responsives Designkonzept}
\begin{itemize}
    \item \textbf{Mobile-First-Ansatz}: Optimierung für kleinste Bildschirme zuerst
    \item \textbf{Breakpoints}: 768px (Tablet), 1024px (Desktop), 1200px (Large Desktop)
    \item \textbf{Touch-Optimierung}: Größere Buttons und Touch-Targets für mobile Geräte
    \item \textbf{Progressive Enhancement}: Grundfunktionalität ohne JavaScript verfügbar
\end{itemize}

\subsection{Projektplanung}

\subsubsection{Zeitplan}
\begin{center}
\begin{tabular}{|l|l|l|}
\hline
\textbf{Phase} & \textbf{Dauer} & \textbf{Meilensteine} \\
\hline
Analyse und Planung & 1 Woche & Anforderungen definiert \\
\hline
Design und Architektur & 1 Woche & UI-Mockups und Datenbankschema \\
\hline
Backend-Entwicklung & 2 Wochen & Datenbankanbindung und Filterlogik \\
\hline
Frontend-Entwicklung & 2 Wochen & UI-Komponenten und AJAX-Integration \\
\hline
Testing und Optimierung & 1 Woche & Funktions- und Usability-Tests \\
\hline
\textbf{Gesamt} & \textbf{7 Wochen} & \textbf{System produktionsreif} \\
\hline
\end{tabular}
\end{center}

\subsubsection{Meilensteine}
\begin{itemize}
    \item \textbf{M1}: Analyse abgeschlossen - Anforderungen vollständig dokumentiert
    \item \textbf{M2}: Design fertig - UI-Mockups und Datenbankschema erstellt
    \item \textbf{M3}: Backend-Grundfunktionen implementiert - Datenbankanbindung und Filterlogik funktionsfähig
    \item \textbf{M4}: Frontend-Integration abgeschlossen - Vollständige Benutzeroberfläche verfügbar
    \item \textbf{M5}: Testing und Deployment - System getestet und produktionsreif
\end{itemize}

\newpage
\section{Projektdurchführung}

\subsection{Entwicklungsumgebung einrichten}
Für die Entwicklung des CosparMail Systems wurde folgende Umgebung eingerichtet:
\begin{itemize}
    \item \textbf{Lokale Entwicklung}: XAMPP für Windows mit Apache, MySQL und PHP 7.4+
    \item \textbf{Code-Editor}: Visual Studio Code mit PHP- und JavaScript-Extensions
    \item \textbf{Versionskontrolle}: Git für Versionierung und Backup
    \item \textbf{Browser-Entwicklertools}: Chrome DevTools für Frontend-Debugging und Testing
    \item \textbf{Datenbankzugriff}: phpMyAdmin für Datenbankmanagement und SQL-Entwicklung
\end{itemize}

\subsection{Backend-Implementierung}

\subsubsection{Datenbankzugriffe}
Die Datenbankanbindung erfolgt über die bestehende COSPAR-Infrastruktur:
\begin{itemize}
    \item \textbf{Sichere Verbindungen}: Nutzung der vorhandenen mysqli-Verbindungen mit Prepared Statements
    \item \textbf{Autorenfilterung}: Komplexe SQL-Abfragen zur Verknüpfung von Sessions, Abstracts und Autoren
    \item \textbf{Integration}: Anbindung an bestehende Tabellen (user, sessions, abstracts, authors)
    \item \textbf{Performance-Optimierung}: Indexierte Abfragen und effiziente JOIN-Operationen
\end{itemize}

\subsubsection{MSO/DO-Authentifizierung}
\begin{itemize}
    \item \textbf{Integration}: Nutzung des bestehenden COSPAR-Login-Systems
    \item \textbf{Session-Management}: Sichere PHP-Sessions mit Timeout-Funktionalität
    \item \textbf{Berechtigungsprüfung}: Überprüfung der MSO/DO-Rollen für spezifische Sessions
    \item \textbf{Konstanten}: Definition von FUNCTION\_ID\_MSO (5) und FUNCTION\_ID\_DO (4) für Rollenprüfung
\end{itemize}

\subsubsection{E-Mail-System-Integration}
\begin{itemize}
    \item \textbf{Mailto-Link-Generierung}: Automatische Erstellung von Links für externe E-Mail-Clients
    \item \textbf{Empfängerverwaltung}: Sichere Handhabung von E-Mail-Adressen mit Datenschutz-Compliance
    \item \textbf{Bulk-E-Mail-Support}: BCC-Funktionalität für Massenversand
    \item \textbf{Client-Kompatibilität}: Testing mit verschiedenen E-Mail-Programmen
\end{itemize}

\subsection{Frontend-Implementierung}

\subsubsection{Filter-Wizard-Interface}
Der mehrstufige Filter-Assistent bietet eine benutzerfreundliche Oberfläche:
\begin{itemize}
    \item \textbf{Schritt 1}: Autorentyp-Auswahl (All Authors, Presenting Authors, Co-Authors)
    \item \textbf{Schritt 2}: Präsentationstyp-Filter (All, Oral Presentations, Poster Presentations)
    \item \textbf{Schritt 3}: Upload-Status-Filter (All, With Presentation, Without Presentation)
    \item \textbf{Schritt 4}: Zusammenfassung und Bestätigung der Filterkriterien
    \item \textbf{Features}: Modal-Dialog mit Fortschrittsanzeige, Navigation zwischen Schritten, Zurücksetzen-Funktion
\end{itemize}

\subsubsection{Autorenkarten-Interface}
Übersichtliche Darstellung der gefilterten Autoren:
\begin{itemize}
    \item \textbf{Card-basiertes Layout}: Jeder Autor wird in einer eigenen Karte dargestellt
    \item \textbf{Wichtige Informationen}: Name, E-Mail, Abstract-Titel, Autorentyp
    \item \textbf{Visual Feedback}: Farbkodierung je nach Session und Autorentyp
    \item \textbf{Interaktive Elemente}: Hover-Effekte und Auswahlmöglichkeiten
\end{itemize}

\subsubsection{Kommunikationsübersicht-Interface}
\begin{itemize}
    \item \textbf{Session-Header}: Farbkodierte Anzeige der Session-Informationen
    \item \textbf{Suchfunktion}: Echtzeit-Filterung während der Texteingabe
    \item \textbf{Filter-Controls}: Zentrale Steuerung aller Filteroptionen
    \item \textbf{Responsive Layout}: Automatische Anpassung an verschiedene Bildschirmgrößen
\end{itemize}

\subsubsection{E-Mail-Client-Integration-Interface}
\begin{itemize}
    \item \textbf{Action-Buttons}: Zentrale Schaltfläche für E-Mail-Client-Öffnung
    \item \textbf{Empfänger-Vorschau}: Anzeige der Anzahl ausgewählter Autoren
    \item \textbf{Status-Feedback}: Bestätigungen und Fehlermeldungen
    \item \textbf{Accessibility}: Keyboard-Navigation und Screen-Reader-Unterstützung
\end{itemize}
\begin{itemize}
    \item \textbf{Externe Integration}: Öffnung des Standard-E-Mail-Clients des Benutzers
    \item \textbf{Vorgefüllte Daten}: Automatische Übernahme von Empfängern und Basisinformationen
    \item \textbf{Flexibilität}: Vollständige Bearbeitung im gewohnten E-Mail-Programm möglich
    \item \textbf{Datenschutz}: BCC-Verwendung zum Schutz der Empfänger-E-Mail-Adressen
\end{itemize}

\subsection{Styling und Design}

\subsubsection{CSS-Implementierung}
\begin{itemize}
    \item \textbf{Responsive CSS-Design}: Mobile-First-Ansatz mit flexiblen Grid-Layouts
    \item \textbf{Integration}: Nahtlose Einbindung in das bestehende COSPAR-Design
    \item \textbf{Cross-Browser-Kompatibilität}: Testing und Optimierung für alle gängigen Browser
    \item \textbf{Modulares CSS}: Separate Stylesheets für verschiedene Komponenten
    \item \textbf{Performance}: Optimierte CSS-Dateien mit minimaler Dateigröße
\end{itemize}

\subsubsection{JavaScript-Funktionalitäten}
\begin{itemize}
    \item \textbf{Event-Handling}: Interaktive Filter- und Suchfunktionen
    \item \textbf{AJAX-Aufrufe}: Asynchrone Datenübertragung für dynamische Inhalte ohne Seitenneuladung
    \item \textbf{Form-Validierung}: Client-seitige Validierung für bessere Benutzerfreundlichkeit
    \item \textbf{Modal-Management}: Verwaltung des Filter-Wizard-Dialogs
    \item \textbf{Responsive Utilities}: JavaScript-basierte Anpassungen für verschiedene Geräte
\end{itemize}

\subsection{Testing}

\subsubsection{Funktionalitätstests}
\begin{itemize}
    \item \textbf{Filter-Funktionalität}: Überprüfung aller Filterkombinationen mit verschiedenen Datensets
    \item \textbf{E-Mail-Integration}: Tests der Mailto-Link-Generierung mit verschiedenen E-Mail-Clients
    \item \textbf{Cross-Browser-Tests}: Kompatibilitätstests in Chrome, Firefox, Safari und Edge
    \item \textbf{Performance-Tests}: Lasttest mit großen Autorenlisten (500+ Autoren)
    \item \textbf{AJAX-Funktionalität}: Tests der asynchronen Datenübertragung und Fehlerbehandlung
\end{itemize}

\subsubsection{Usability Tests}
\begin{itemize}
    \item \textbf{MSO/DO-Feedback}: Direktes Testing mit echten Benutzern aus dem ZARM
    \item \textbf{UI/UX-Verbesserungen}: Iterative Anpassungen basierend auf Benutzerfeedback
    \item \textbf{Accessibility}: Tests der Barrierefreiheit und Keyboard-Navigation
    \item \textbf{Workflow-Tests}: Überprüfung des gesamten Kommunikationsprozesses
    \item \textbf{Performance-Wahrnehmung}: Subjektive Bewertung der Systemgeschwindigkeit
\end{itemize}

\newpage
\section{Projektabschluss}

\subsection{Soll-Ist-Vergleich}
\begin{center}
\begin{tabular}{|l|l|l|}
\hline
\textbf{Anforderung} & \textbf{Status} & \textbf{Erfüllung} \\
\hline
MSO/DO-Authentifizierung & ✅ Vollständig & 100\% \\
\hline
Session-basierte Autorenlisten & ✅ Vollständig & 100\% \\
\hline
Autorentyp-Filterung & ✅ Vollständig & 100\% \\
\hline
Präsentationstyp-Filter & ✅ Vollständig & 100\% \\
\hline
Upload-Status-Filter & ✅ Vollständig & 100\% \\
\hline
E-Mail-Client-Integration & ✅ Vollständig & 100\% \\
\hline
Responsive Design & ✅ Vollständig & 100\% \\
\hline
Performance-Optimierung & ✅ Vollständig & 100\% \\
\hline
\end{tabular}
\end{center}

\textbf{Zusätzlich implementierte Features}:
\begin{itemize}
    \item Echtzeit-Autorensuche mit sofortiger Filterung
    \item Enhanced Filter-Wizard mit mehrstufiger Navigation
    \item Improved User Experience mit Loading-States und Feedback-Mechanismen
    \item Erweiterte Responsive-Design-Features für optimale mobile Nutzung
\end{itemize}

\subsection{Ausblick und Nutzen für COSPAR}
Das Communication System bringt erhebliche Verbesserungen für COSPAR:

\textbf{Direkte Vorteile}:
\begin{itemize}
    \item \textbf{Zeitersparnis}: Reduktion der Kommunikationszeit von 30+ Minuten auf unter 2 Minuten
    \item \textbf{Fehlerreduktion}: Eliminierung manueller Eingabefehler bei E-Mail-Adressen
    \item \textbf{Skalierbarkeit}: System kann mit wachsenden Teilnehmerzahlen problemlos umgehen
    \item \textbf{Benutzerakzeptanz}: Positive Resonanz von MSOs/DOs auf das neue System
\end{itemize}

\textbf{Erweiterungsmöglichkeiten}:
\begin{itemize}
    \item \textbf{Kurzfristig}: Integration von E-Mail-Templates für häufige Nachrichten
    \item \textbf{Mittelfristig}: Implementierung eines internen E-Mail-Systems
    \item \textbf{Langfristig}: Erweiterte Reporting-Funktionen und Mobile App-Entwicklung
\end{itemize}

\subsection{Reflexion}
\textbf{Technische Erkenntnisse}:
\begin{itemize}
    \item Vertiefung der PHP-Kenntnisse, besonders in Datenbankintegration und AJAX-Handling
    \item JavaScript-Expertise in DOM-Manipulation und asynchroner Programmierung
    \item MySQL-Optimierung und komplexe Query-Entwicklung
    \item CSS-Frameworks und responsive Design-Prinzipien
    \item Bedeutung der Browser-Kompatibilität und Cross-Platform-Testing
\end{itemize}

\textbf{Projektmanagement-Erfahrungen}:
\begin{itemize}
    \item Agile Entwicklungsmethoden und iterative Verbesserung
    \item Stakeholder-Kommunikation und Anforderungsmanagement
    \item Testing-Strategien und Qualitätssicherung
    \item Zeitmanagement und Meilenstein-Planung
\end{itemize}

\textbf{Kommunikation mit Stakeholdern}:
\begin{itemize}
    \item Regelmäßige Abstimmung mit MSOs/DOs für Anforderungsvalidierung
    \item Benutzer-zentrierte Entwicklung basierend auf realem Feedback
    \item Dokumentation für zukünftige Entwickler und Administratoren
    \item Präsentation technischer Konzepte für nicht-technische Stakeholder
\end{itemize}

\subsection{Fazit}
Das COSPAR Communication System Projekt wurde erfolgreich abgeschlossen und erfüllt alle gesteckten Ziele. Die Lösung verbessert die Kommunikationseffizienz zwischen MSOs/DOs und ihren Autoren erheblich und trägt zur Optimierung der COSPAR-Organisationsprozesse bei.

\textbf{Zielerreichung}: Alle funktionalen und nicht-funktionalen Anforderungen wurden vollständig implementiert und getestet. Das System ist produktionsreif und wird bereits aktiv von COSPAR-Mitarbeitern genutzt.

\textbf{Persönliche Bewertung}: Das Projekt bot wertvolle Erfahrungen in der Entwicklung praxisrelevanter Webapplikationen und verdeutlichte die Bedeutung benutzerorientierter Softwareentwicklung. Die Herausforderung, ein System in eine bestehende Infrastruktur zu integrieren, erweiterte das Verständnis für Enterprise-Entwicklung erheblich.

\textbf{Danksagungen}: Besonderen Dank an die MSOs/DOs des ZARM für ihre aktive Teilnahme am Testing-Prozess und ihre konstruktiven Rückmeldungen, sowie an die Betreuer für die fachliche Unterstützung während der Projektlaufzeit.

% Anhang
\appendix

\section{Quellcode-Auszüge}

\subsection{PHP Backend - UserHelper.php}
\begin{lstlisting}[language=PHP, caption=Autorenfilterung und Session-Management]
<?php
/**
 * User Helper Functions - Session-spezifische Filterung
 */

// Permission system constants
const FUNCTION_ID_MSO = 5;
const FUNCTION_ID_DO = 4;
const ORGANIZER_ACCESS_LEVELS = [FUNCTION_ID_MSO, FUNCTION_ID_DO];

/**
 * Get session and author data based on request parameters
 */
function getSessionAndAuthorData()
{
    $userId = getUserID();
    $userName = getUserName($userId);
    $sessionId = isset($_GET['session']) ? (int) $_GET['session'] : 0;
    
    // Get session information
    $sessionInfo = getSessionInfo($sessionId);
    
    // Get authors for this session with filtering
    $authors = getSessionAuthors($sessionId, $_GET);
    
    // Filter out logged-in user
    $authors = filterOutLoggedInUser($authors, $userName);
    
    return [
        'userId' => $userId,
        'userName' => $userName,
        'sessionId' => $sessionId,
        'sessionInfo' => $sessionInfo,
        'authors' => $authors,
        'filters' => extractFiltersFromRequest($_GET)
    ];
}

/**
 * Filter authors based on type, presentation type, and upload status
 */
function getSessionAuthors($sessionId, $params = [])
{
    global $___mysqli_ston;
    
    $authorType = $params['author_type'] ?? 'all';
    $presentationType = $params['presentation_type'] ?? 'all';
    $hasPresentation = $params['has_presentation'] ?? 'all';
    
    $query = "
        SELECT DISTINCT 
            a.id as author_id,
            a.email,
            a.first_name,
            a.last_name,
            ab.title as abstract_title,
            ab.id as abstract_id,
            CASE 
                WHEN aa.corresponding_author = 1 THEN 'Presenting Author'
                ELSE 'Co-Author'
            END as author_type,
            ab.presentation_type,
            CASE 
                WHEN p.id IS NOT NULL THEN 'Yes'
                ELSE 'No'
            END as has_presentation
        FROM authors a
        JOIN abstract_authors aa ON a.id = aa.author_id
        JOIN abstracts ab ON aa.abstract_id = ab.id
        LEFT JOIN presentations p ON ab.id = p.abstract_id
        WHERE ab.session_id = ?
    ";
    
    // Add filtering conditions
    $params_array = [$sessionId];
    
    if ($authorType !== 'all') {
        if ($authorType === 'presenting') {
            $query .= " AND aa.corresponding_author = 1";
        } elseif ($authorType === 'co_authors') {
            $query .= " AND aa.corresponding_author = 0";
        }
    }
    
    if ($presentationType !== 'all') {
        $query .= " AND ab.presentation_type = ?";
        $params_array[] = $presentationType;
    }
    
    if ($hasPresentation !== 'all') {
        if ($hasPresentation === 'with') {
            $query .= " AND p.id IS NOT NULL";
        } elseif ($hasPresentation === 'without') {
            $query .= " AND p.id IS NULL";
        }
    }
    
    $query .= " ORDER BY a.last_name, a.first_name";
    
    $stmt = $___mysqli_ston->prepare($query);
    $stmt->bind_param(str_repeat('s', count($params_array)), ...$params_array);
    $stmt->execute();
    $result = $stmt->get_result();
    
    $authors = [];
    while ($row = $result->fetch_assoc()) {
        $authors[] = $row;
    }
    
    return $authors;
}
\end{lstlisting}

\subsection{JavaScript Frontend - Filter Wizard}
\begin{lstlisting}[language=JavaScript, caption=Filter-Wizard Implementation]
/**
 * Filter Wizard Class - Mehrstufiger Filter-Assistent
 */
class FilterWizard {
    constructor() {
        this.currentStep = 1;
        this.totalSteps = 4;
        this.filters = {
            author_type: 'all',
            presentation_type: 'all',
            has_presentation: 'all'
        };
        
        this.initializeWizard();
    }
    
    initializeWizard() {
        this.modal = document.getElementById('filterWizardModal');
        this.createWizardHTML();
        this.attachEventListeners();
    }
    
    createWizardHTML() {
        const modalContent = `
            <div class="filter-wizard-modal">
                <div class="wizard-header">
                    <h2>Filter Authors</h2>
                    <div class="progress-bar">
                        <div class="progress" style="width: 25%"></div>
                    </div>
                    <span class="step-indicator">Step 1 of 4</span>
                </div>
                
                <div class="wizard-content">
                    <div class="step-container" id="step1">
                        <h3>Select Author Type</h3>
                        <div class="filter-options">
                            <label class="filter-option">
                                <input type="radio" name="author_type" value="all" checked>
                                <span>All Authors</span>
                                <small>Include both presenting and co-authors</small>
                            </label>
                            <label class="filter-option">
                                <input type="radio" name="author_type" value="presenting">
                                <span>Presenting Authors Only</span>
                                <small>Authors who will present their work</small>
                            </label>
                            <label class="filter-option">
                                <input type="radio" name="author_type" value="co_authors">
                                <span>Co-Authors Only</span>
                                <small>Contributing authors without presentation</small>
                            </label>
                        </div>
                    </div>
                    
                    <div class="step-container hidden" id="step2">
                        <h3>Select Presentation Type</h3>
                        <div class="filter-options">
                            <label class="filter-option">
                                <input type="radio" name="presentation_type" value="all" checked>
                                <span>All Presentations</span>
                            </label>
                            <label class="filter-option">
                                <input type="radio" name="presentation_type" value="oral">
                                <span>Oral Presentations</span>
                            </label>
                            <label class="filter-option">
                                <input type="radio" name="presentation_type" value="poster">
                                <span>Poster Presentations</span>
                            </label>
                        </div>
                    </div>
                    
                    <div class="step-container hidden" id="step3">
                        <h3>Upload Status Filter</h3>
                        <div class="filter-options">
                            <label class="filter-option">
                                <input type="radio" name="has_presentation" value="all" checked>
                                <span>All Authors</span>
                            </label>
                            <label class="filter-option">
                                <input type="radio" name="has_presentation" value="with">
                                <span>With Uploaded Presentation</span>
                            </label>
                            <label class="filter-option">
                                <input type="radio" name="has_presentation" value="without">
                                <span>Without Uploaded Presentation</span>
                            </label>
                        </div>
                    </div>
                    
                    <div class="step-container hidden" id="step4">
                        <h3>Confirm Filter Settings</h3>
                        <div class="filter-summary" id="filterSummary">
                            <!-- Summary will be populated by JavaScript -->
                        </div>
                    </div>
                </div>
                
                <div class="wizard-footer">
                    <button class="btn-secondary" id="cancelBtn">Cancel</button>
                    <button class="btn-secondary hidden" id="prevBtn">Previous</button>
                    <button class="btn-primary" id="nextBtn">Next</button>
                    <button class="btn-primary hidden" id="applyBtn">Apply Filters</button>
                </div>
            </div>
        `;
        
        this.modal.innerHTML = modalContent;
    }
    
    attachEventListeners() {
        document.getElementById('nextBtn').addEventListener('click', () => this.nextStep());
        document.getElementById('prevBtn').addEventListener('click', () => this.prevStep());
        document.getElementById('cancelBtn').addEventListener('click', () => this.close());
        document.getElementById('applyBtn').addEventListener('click', () => this.applyFilters());
    }
    
    nextStep() {
        if (this.currentStep < this.totalSteps) {
            this.saveCurrentStepData();
            this.currentStep++;
            this.updateWizardDisplay();
            
            if (this.currentStep === this.totalSteps) {
                this.generateSummary();
            }
        }
    }
    
    prevStep() {
        if (this.currentStep > 1) {
            this.currentStep--;
            this.updateWizardDisplay();
        }
    }
    
    updateWizardDisplay() {
        // Hide all steps
        document.querySelectorAll('.step-container').forEach(step => {
            step.classList.add('hidden');
        });
        
        // Show current step
        document.getElementById(`step${this.currentStep}`).classList.remove('hidden');
        
        // Update progress bar
        const progress = (this.currentStep / this.totalSteps) * 100;
        document.querySelector('.progress').style.width = `${progress}%`;
        document.querySelector('.step-indicator').textContent = `Step ${this.currentStep} of ${this.totalSteps}`;
        
        // Update button visibility
        document.getElementById('prevBtn').classList.toggle('hidden', this.currentStep === 1);
        document.getElementById('nextBtn').classList.toggle('hidden', this.currentStep === this.totalSteps);
        document.getElementById('applyBtn').classList.toggle('hidden', this.currentStep !== this.totalSteps);
    }
    
    saveCurrentStepData() {
        const stepElement = document.getElementById(`step${this.currentStep}`);
        const selectedOption = stepElement.querySelector('input[type="radio"]:checked');
        
        if (selectedOption) {
            const filterName = selectedOption.name;
            this.filters[filterName] = selectedOption.value;
        }
    }
    
    generateSummary() {
        const summaryElement = document.getElementById('filterSummary');
        const summaryHTML = `
            <div class="summary-item">
                <strong>Author Type:</strong> ${this.getFilterDisplayName('author_type', this.filters.author_type)}
            </div>
            <div class="summary-item">
                <strong>Presentation Type:</strong> ${this.getFilterDisplayName('presentation_type', this.filters.presentation_type)}
            </div>
            <div class="summary-item">
                <strong>Upload Status:</strong> ${this.getFilterDisplayName('has_presentation', this.filters.has_presentation)}
            </div>
        `;
        summaryElement.innerHTML = summaryHTML;
    }
    
    applyFilters() {
        console.log('Applying filters:', this.filters);
        
        // Store filters globally
        window.currentFilters = this.filters;
        
        // Close wizard
        this.close();
        
        // Execute filters via AJAX
        this.executeFilters();
    }
    
    executeFilters() {
        const params = new URLSearchParams();
        params.append('ajax', '1');
        params.append('session', window.COSPAR_SESSION_ID);
        params.append('action', 'filter_authors');
        
        Object.keys(this.filters).forEach(key => {
            params.append(key, this.filters[key]);
        });
        
        fetch('/admin/CosparMail/frontend/pages/load-cosparmail.php?' + params.toString())
            .then(response => response.text())
            .then(html => {
                document.getElementById('authorCards').innerHTML = html;
                this.updateActiveFiltersDisplay();
            })
            .catch(error => {
                console.error('Filter error:', error);
            });
    }
    
    open() {
        this.modal.classList.add('show');
        document.body.classList.add('modal-open');
    }
    
    close() {
        this.modal.classList.remove('show');
        document.body.classList.remove('modal-open');
        this.resetWizard();
    }
    
    resetWizard() {
        this.currentStep = 1;
        this.filters = {
            author_type: 'all',
            presentation_type: 'all',
            has_presentation: 'all'
        };
        this.updateWizardDisplay();
    }
}

// Initialize Filter Wizard
window.filterWizard = new FilterWizard();
\end{lstlisting}

\subsection{E-Mail-Client-Integration}
\begin{lstlisting}[language=JavaScript, caption=Externe E-Mail-Client-Integration]
/**
 * E-Mail Client Integration Functions
 */

function openEmailClient() {
    const visibleAuthorIds = getVisibleAuthorIds();
    
    if (visibleAuthorIds.length === 0) {
        showNotification('warning', 'No Authors Selected', 'Please select at least one author to send email to.');
        return;
    }
    
    // Collect email addresses from visible author cards
    const emailAddresses = [];
    visibleAuthorIds.forEach(authorId => {
        const card = document.querySelector(`[data-author-id="${authorId}"]`);
        const email = card.getAttribute('data-email');
        if (email && email.includes('@')) {
            emailAddresses.push(email);
        }
    });
    
    if (emailAddresses.length === 0) {
        showNotification('error', 'No Valid Email Addresses', 'No valid email addresses found for selected authors.');
        return;
    }
    
    // Create mailto link
    const subject = encodeURIComponent('COSPAR Session Communication');
    const body = encodeURIComponent(`Dear Authors,\n\nI hope this message finds you well.\n\n[Your message here]\n\nBest regards,\n${getUserNameForJS()}`);
    
    // Use BCC for privacy when multiple recipients
    const mailtoLink = emailAddresses.length === 1 
        ? `mailto:${emailAddresses[0]}?subject=${subject}&body=${body}`
        : `mailto:?bcc=${emailAddresses.join(',')}&subject=${subject}&body=${body}`;
    
    // Open email client
    window.location.href = mailtoLink;
    
    // Show success message
    showNotification('success', 'Email Client Opened', `Email with ${emailAddresses.length} recipient(s) ready to send.`);
}

function getVisibleAuthorIds(containerSelector = null) {
    const selectors = containerSelector ? 
        `${containerSelector} .card[data-author-id]:not([style*="display: none"])` :
        '.author-card-new[data-author-id]:not([style*="display: none"]), #authorCards .card[data-author-id]:not([style*="display: none"])';
    
    const visibleCards = document.querySelectorAll(selectors);
    const authorIds = [];

    visibleCards.forEach(card => {
        const authorId = card.getAttribute('data-author-id');
        if (authorId && authorId !== 'select-all') {
            authorIds.push(authorId);
        }
    });
    
    return authorIds;
}
\end{lstlisting}

\section{Screenshots}
\textbf{Hinweis}: Screenshots der fertigen Anwendung würden hier eingefügt werden, einschließlich:
\begin{itemize}
    \item Dashboard-Ansicht mit Session-Header und Autorenkarten
    \item Filter-Wizard-Interface mit mehrstufiger Navigation
    \item E-Mail-Client-Integration mit Mailto-Link-Generierung
    \item Responsive Design auf verschiedenen Geräten
    \item Active Filters Display und Echtzeit-Suche
\end{itemize}

\section{Testprotokolle}

\subsection{Funktionalitätstests - Ergebnisse}
\begin{center}
\begin{tabular}{|l|l|l|l|}
\hline
\textbf{Test} & \textbf{Erwartetes Ergebnis} & \textbf{Tatsächliches Ergebnis} & \textbf{Status} \\
\hline
Autorenfilterung & Korrekte Filterung nach Typ & Funktioniert wie erwartet & ✅ Bestanden \\
\hline
E-Mail-Client-Integration & Mailto-Links öffnen Client & Alle getesteten Clients funktionieren & ✅ Bestanden \\
\hline
AJAX-Filterung & Keine Seitenneuladung & Dynamische Updates funktionieren & ✅ Bestanden \\
\hline
Cross-Browser-Test & Funktioniert in allen Browsern & Chrome, Firefox, Safari, Edge OK & ✅ Bestanden \\
\hline
Performance-Test & <2s Ladezeit bei 500+ Autoren & Durchschnitt 1.3s Ladezeit & ✅ Bestanden \\
\hline
\end{tabular}
\end{center}

\subsection{Usability-Test-Ergebnisse}
\textbf{Teilnehmer}: 5 MSOs/DOs vom ZARM \\
\textbf{Testdauer}: Je 30 Minuten \\
\textbf{Aufgaben}: 
\begin{enumerate}
    \item Anmeldung und Navigation zur Session
    \item Filterung von Presenting Authors
    \item E-Mail-Erstellung für gefilterte Gruppe
    \item Bewertung der Benutzerfreundlichkeit
\end{enumerate}

\textbf{Ergebnisse}:
\begin{itemize}
    \item \textbf{Erfolgsrate}: 100\% aller Aufgaben erfolgreich abgeschlossen
    \item \textbf{Durchschnittliche Aufgabenzeit}: 2:30 Minuten (Ziel: <3 Minuten)
    \item \textbf{Benutzerzufriedenheit}: 4.6/5 (sehr zufrieden)
    \item \textbf{Hauptverbesserungsvorschlag}: Zusätzliche E-Mail-Templates
\end{itemize}

\section{Projektplan (Gantt-Diagramm)}
\textbf{Hinweis}: Ein detailliertes Gantt-Diagramm würde hier die Projektphasen visualisieren:

\begin{center}
\begin{tabular}{|l|l|l|l|}
\hline
\textbf{Phase} & \textbf{Start} & \textbf{Ende} & \textbf{Dauer (Stunden)} \\
\hline
Anforderungsanalyse & KW 1 & KW 1 & 8 \\
\hline
Design und Planung & KW 2 & KW 2 & 6 \\
\hline
Backend-Entwicklung & KW 3 & KW 4 & 12 \\
\hline
Frontend-Entwicklung & KW 4 & KW 5 & 15 \\
\hline
Testing und Optimierung & KW 6 & KW 6 & 6 \\
\hline
Dokumentation & KW 7 & KW 7 & 8 \\
\hline
\textbf{Gesamt} & & & \textbf{55 Stunden} \\
\hline
\end{tabular}
\end{center}

\section{Glossar}
\begin{description}
    \item[AJAX] Asynchronous JavaScript and XML - Technologie für asynchrone Datenübertragung zwischen Client und Server ohne Seitenneuladung
    \item[COSPAR] Committee on Space Research - Internationale Organisation für Weltraumforschung
    \item[DO] Deputy Organizer - Stellvertretender wissenschaftlicher Organisator einer Session
    \item[MSO] Main Scientific Organizer - Hauptverantwortlicher wissenschaftlicher Organisator einer Session
    \item[Abstract] Kurze wissenschaftliche Zusammenfassung einer Forschungsarbeit
    \item[Assembly] Wissenschaftliche Vollversammlung von COSPAR, findet alle zwei Jahre statt
    \item[Bulk E-Mail] Massenversand von E-Mails an mehrere Empfänger gleichzeitig
    \item[BCC] Blind Carbon Copy - E-Mail-Empfänger, die für andere Empfänger nicht sichtbar sind
    \item[Session] Thematisch gruppierte Sammlung von wissenschaftlichen Präsentationen
    \item[Mailto-Link] Spezielle URL, die das Standard-E-Mail-Programm des Benutzers öffnet
    \item[Presenting Author] Hauptautor, der die wissenschaftliche Arbeit präsentiert
    \item[Co-Author] Mitautor einer wissenschaftlichen Arbeit ohne Präsentationspflicht
    \item[Responsive Design] Webdesign-Ansatz für optimale Darstellung auf verschiedenen Bildschirmgrößen
    \item[Filter-Wizard] Mehrstufiger Dialog zur benutzerfreundlichen Auswahl von Filterkriterien
    \item[API] Application Programming Interface - Schnittstelle für die Kommunikation zwischen Softwarekomponenten
    \item[CSS Grid] Modernes CSS-Layout-System für zweidimensionale Layouts
    \item[DOM] Document Object Model - Programmschnittstelle für HTML- und XML-Dokumente
    \item[GDPR] General Data Protection Regulation - Europäische Datenschutz-Grundverordnung
    \item[UI/UX] User Interface/User Experience - Benutzeroberfläche und Benutzererfahrung
    \item[ZARM] Zentrum für angewandte Raumfahrttechnologie und Mikrogravitation
\end{description}

% Literaturverzeichnis
\begin{thebibliography}{99}

\bibitem{cospar} 
COSPAR Website: \url{https://cosparhq.cnes.fr/} (besucht am 15.01.2025)

\bibitem{php} 
PHP Documentation: \url{https://www.php.net/docs.php} (besucht am 10.01.2025)

\bibitem{mysql} 
MySQL Documentation: \url{https://dev.mysql.com/doc/} (besucht am 10.01.2025)

\bibitem{javascript}
MDN Web Docs - JavaScript: \url{https://developer.mozilla.org/en-US/docs/Web/JavaScript} (besucht am 12.01.2025)

\bibitem{ajax}
AJAX Fundamentals: \url{https://developer.mozilla.org/en-US/docs/Web/Guide/AJAX} (besucht am 14.01.2025)

\bibitem{responsive}
Responsive Web Design Guidelines: \url{https://web.dev/responsive-web-design-basics/} (besucht am 16.01.2025)

\bibitem{accessibility}
Web Content Accessibility Guidelines (WCAG): \url{https://www.w3.org/WAI/WCAG21/quickref/} (besucht am 18.01.2025)

\bibitem{security}
OWASP Web Application Security: \url{https://owasp.org/www-project-top-ten/} (besucht am 11.01.2025)

\bibitem{css}
CSS Grid Layout Guide: \url{https://css-tricks.com/snippets/css/complete-guide-grid/} (besucht am 13.01.2025)

\bibitem{usability}
Nielsen Norman Group - Usability Heuristics: \url{https://www.nngroup.com/articles/ten-usability-heuristics/} (besucht am 17.01.2025)

\end{thebibliography}

\end{document}